\documentclass[10pt]{article}
\usepackage[margin=12mm]{geometry}
\usepackage{graphicx}
\usepackage{caption}
\pagestyle{empty}
\captionsetup{font=small,labelfont=bf}

\begin{document}
\begin{figure}[p]
\centering
\includegraphics[width=0.95\textwidth]{fig6A_ce_retention_heatmap.pdf}\\[1mm]
\includegraphics[width=0.95\textwidth]{fig6B_ce_within_type_fronts.pdf}\\[1mm]
\includegraphics[width=0.95\textwidth]{fig6C_ce_type_restricted_aggregate.pdf}

\caption{\textbf{Terminal cell types retain distinct lineage structure,
while type-restricted optimization recovers only part of the global
travel--cell-state trade-off.}
%
\textbf{(A)} Natural-edge retention for each terminal cell type across the
global Pareto front. Cell types individually represented by four or fewer
cells are pooled into "other" type. Position is normalized front arc length in the
endpoint-normalized coordinates defined in Figure 5A. Profiles are resampled
to uniform intervals and locally averaged for display only; the assignments
and marked keypoint are unchanged. Brighter blue denotes greater retention;
the scale is anchored at zero and spans the displayed maximum without
clipping. The blue diamond and dashed line mark the Pareto assignment with
maximum natural-edge retention.
%
\textbf{(B)} Pareto fronts obtained by allowing terminal assignments only
within each displayed type. For comparability, changes in both objectives are
expressed per cell in standard-deviation units from the embryo-wide
first-cousin null; this null sets the common scale only and does not constrain
the assignments. The cross marks the natural assignment;
teal circles and vermillion squares mark the travel and cell-state optima,
respectively.
%
\textbf{(C)} Aggregate front when assignments are prohibited across all
terminal cell-type boundaries (orange), compared with the unrestricted
front (charcoal), using the same global scaling and natural-lineage origin.
Type-preserving assignments form a subset of the unrestricted feasible set.
At the corresponding single-objective optima, this restriction retains an
additional 3.10 standard deviations of travel distance and 15.07 standard
deviations of cell-state distance relative to unrestricted optimization.}
\label{fig:ce_cell_types}
\end{figure}
\end{document}
